\section{Ethical Considerations}
The reported new QUIC experiments use local processes, demo identities,
simulated customer approval, and fixed or simulated business effects. They
make no calls to unsuspecting contact-center workers and do not collect
private production conversations. Pindrop's report is a cited public source;
its underlying call data were not obtained or independently reanalyzed.
The new verification runs require no paid LLM or carrier API calls.

Real-carrier or acoustic experiments must use consented participants and
controlled receiving accounts, comply with recording rules, avoid emergency
and third-party numbers, and prevent real financial or physical effects.
Prompt-injection testing should use synthetic secrets and sandboxed tools.
Human studies need appropriate institutional review and accessible disclosure
of modem audibility, failures, and withdrawal options. We do not claim such
approval or studies have been completed. Artifacts must not release private
keys, account tokens, production audio, or identifying customer metadata.

\section{Open Science and Artifact Access}
The research draft accompanies the source repository. The implementation
baseline is commit \implementationcommit{}, and the sanitized measurement
record is \path{docs/quic-verification.json}. This file records environment
versions, single-run timings, authorization outcomes, and fault counters.
The paper's result table is regenerated by
\path{paper/scripts/extract_results.py}. No new statistical study or formal
verification artifact is represented as complete.

Reproduction requires Python 3.11+, an OpenSSL executable for certificate
demos, and the optional QUIC, audio, authority, WebRTC, Twilio, and PAP
dependencies in the lockfile. Local emulation needs neither an OpenAI API key
nor a Twilio account. A complete run is:
\begin{verbatim}
uv sync --locked --extra quic --extra webrtc \
  --extra authority --extra twilio --extra pap
uv run --no-sync pytest -q
uv run --no-sync openavt sdk-demo \
  --protocol quic --loss --output output/repro
\end{verbatim}
The companion README lists provider, PAP, and FSK commands, the original
reports, template provenance, and the evidence still required before
submission. An external-review artifact must be deposited at a stable,
anonymous location before a double-blind submission. This working tree is
not that anonymous deposit; no reviewer access link or submission identifier
has been invented.

\section{Evidence Boundaries}
The table in the main text contains four single-call observations. The FSK
row has no injected faults; the lossy rows use the fixed waveform impairment
schedule, not a random-loss distribution. Passing all eight SDK probes means
one allowed handler and seven denied requests in those fixtures. It is not a
measured false-acceptance rate for arbitrary agents. Exactly-once confirmation
is a property of the local PAP reference operation fixture, not of QUIC alone.
Historical TLS measurements and current QUIC measurements differ in environment
and are not plotted as a causal performance comparison.
